\section{Extended Related Work}
\label{app:extended_related_work}
There have been a lot of social science works that have done agent-based modeling to study human interactions, spanning across various domains such as economics, phychology, and education \citep{Sawyer_2005, Rosé2008, DEGUCHI1995269}.
Prior simulation environments have played a pivotal role in constructing theories and generating hypotheses in these fields. 
However, they frequently constrain agents' communicative capacities to artificial languages and present a highly reductionist view of simulated human behavior \citep{Gilbert2005-cp, Tesfatsion2006-is, Huang2014-bk, kovač2021socialai, urbanek-etal-2019-learning}.
LLMs provide a more flexible and expressive way to model human behavior. Here, we include a more detailed discussion of the recent works investigating LLMs for simulating human social interactions.
There are works that focus on investigating the fidelity of LLMs in keeping the designated persona and experiences of the characters \citep{shao-etal-2023-character, Jiang2023PersonaLLMIT}.
There are works that simulate human social interactions focusing on certain aspects such as competition, collaboration, negotiation, deception, problem-sovling and etc., \citep{zhang2024exploring, zhao2023competeai, Liu2023ImprovingIC, michael2023debate, Rasal2024LLMHM, hubinger2024sleeper, bianchi2024llms, Xie2024CanLL, Jiang2023PersonaLLMIT}.
As LLMs are becoming more and more popular in simulating human social interactions, there are also works that focus on investigating the potential issues and challenges of using LLMs in social simulations, such as stereotypes and reporting issues \citep{cheng-etal-2023-compost, zhou2024real}.


